\chapter{Execution Plan\sv{Genomförandeplan}}
\label{chap:plan}
\begin{instructions}
\begin{onlymidterm}
This chapter is a refined version of the execution plan presented in the Project Plan. The timetable (Gantt chart) must be updated.

\end{onlymidterm}
Tell the reader how you plan to organise your work. Include a timetable (Gantt chart) showing milestones and planned deadlines that include:
\begin{itemize}
    \item deliverables and deadlines to supervisor ahead of examinations,
    \item deliverables and deadlines to examiner ahead of examinations,
    \item deliverables and deadlines for informing yourselves for the technical background,
    \item deadlines for practical work.
\end{itemize}
\begin{onlyplan}

Figure~\ref{fig:gantt-plan} is an example of a Gantt chart for a 20-week project, made with the \texttt{pgfgantt} package. Bars are activities and diamonds are milestones (deadlines). Copy the code and change the activities, weeks and milestones to fit your own project. Week 1 is the first week of the course, not the calendar week.

\begin{figure}[htbp]
  \centering
  \begin{ganttchart}[hhgantt]{1}{20}
    \gantttitle{Week}{20} \\
    \gantttitlelist{1,...,20}{1} \\
    \ganttgroup{Planning}{1}{3} \\
    \ganttbar{Project plan}{1}{3} \\
    \ganttmilestone{Project plan to examiner}{3} \\
    \ganttgroup{Technical background}{2}{9} \\
    \ganttbar{Literature study}{2}{5} \\
    \ganttbar{Write technical background}{4}{9} \\
    \ganttgroup{Practical work}{4}{12} \\
    \ganttbar{Requirements and design}{4}{6} \\
    \ganttbar{First prototype}{6}{8} \\
    \ganttmilestone{Prototype working}{8} \\
    \ganttbar{Implementation}{8}{12} \\
    \ganttbar{Tests and evaluation}{10}{12} \\
    \ganttgroup{Reporting}{8}{20} \\
    \ganttbar{Half-time report}{8}{10} \\
    \ganttmilestone{Half-time report to supervisor}{9} \\
    \ganttmilestone{Half-time report to examiner}{10} \\
    \ganttbar{Write thesis}{11}{18} \\
    \ganttmilestone{Draft to supervisor}{16} \\
    \ganttmilestone{Thesis to examiner}{18} \\
    \ganttmilestone{UtExpo}{19} \\
    \ganttmilestone{Presentation and opposition}{20}
  \end{ganttchart}
  \caption{Example of a Gantt chart for the project plan}
  \label{fig:gantt-plan}
\end{figure}
\end{onlyplan}
\begin{onlymidterm}

Figure~\ref{fig:gantt-midterm} is the Gantt chart from the Project Plan, updated for the half-time report in week 10 (red line). The filled part of each bar shows how much of the activity is done, and the white part what remains. Where the plan has changed, a thin grey bar under the activity shows the original plan. Keep the final deadlines unless you have agreed on new ones with your supervisor and examiner.

\begin{figure}[htbp]
  \centering
  \begin{ganttchart}[hhgantt, today=10, today label={Now}]{1}{20}
    \gantttitle{Week}{20} \\
    \gantttitlelist{1,...,20}{1} \\
    \ganttgroup{Planning}{1}{3} \\
    \ganttbar[progress=100]{Project plan}{1}{3} \\
    \ganttmilestone{Project plan to examiner}{3} \\
    \ganttgroup{Technical background}{2}{10} \\
    \ganttbar[progress=100]{Literature study}{2}{6}
      \ganttbar[hhplanned]{}{2}{5} \\
    \ganttbar[progress=80]{Write technical background}{4}{10}
      \ganttbar[hhplanned]{}{4}{9} \\
    \ganttgroup{Practical work}{4}{14} \\
    \ganttbar[progress=100]{Requirements and design}{4}{6} \\
    \ganttbar[progress=100]{First prototype}{6}{8} \\
    \ganttmilestone{Prototype working}{8} \\
    \ganttbar[progress=0]{Wait for new hardware (new)}{9}{11} \\
    \ganttbar[progress=40]{Implementation}{8}{13}
      \ganttbar[hhplanned]{}{8}{12} \\
    \ganttbar[progress=0]{Tests and evaluation}{12}{14}
      \ganttbar[hhplanned]{}{10}{12} \\
    \ganttgroup{Reporting}{8}{20} \\
    \ganttbar[progress=100]{Half-time report}{8}{10} \\
    \ganttmilestone{Half-time report to supervisor}{9} \\
    \ganttmilestone{Half-time report to examiner}{10} \\
    \ganttbar[progress=0]{Write thesis}{11}{18} \\
    \ganttmilestone{Draft to supervisor}{16} \\
    \ganttmilestone{Thesis to examiner}{18} \\
    \ganttmilestone{UtExpo}{19} \\
    \ganttmilestone{Presentation and opposition}{20}
  \end{ganttchart}
  \caption{Example of an updated Gantt chart for the half-time report}
  \label{fig:gantt-midterm}
\end{figure}
\end{onlymidterm}
\end{instructions}

\begin{onlymidterm}
\section{Follow-up of the Project Plan\sv{Uppföljning av projektplanen}}
\label{sec:planfollowup}
\begin{instructions}
Tell the reader how the work has progressed compared to the Project Plan. If the project plan has changed, this should be stated, for example, if you have changed the requirements specification (Section~\ref{sec:req}) or the project boundaries. Explain why the changes were made and how they affect the remaining work.

Summarise the milestones from the Project Plan in a table like Table~\ref{tab:followup}.

\begin{table}[ht]
  \centering
  \caption{Follow-up of the milestones in the Project Plan}
  \label{tab:followup}
  \begin{tabular}{llll}
    \toprule
    Milestone & Planned & Done / new date & Comment \\
    \midrule
    Literature study done & Week 5 & Week 6 & Took longer than planned \\
    First prototype & Week 8 & Week 8 & \\
    Tests finished & Week 12 & Week 14 (new) & Hardware delivered late \\
    \bottomrule
  \end{tabular}
\end{table}
\end{instructions}
\end{onlymidterm}

\begin{onlyplan}
\section{Chosen area of the Technical background\sv{Valt område för teknisk bakgrund}}
\label{sec:chosentechbg}
\begin{instructions}
In which area should you describe your technical background of the project? Give a brief overview of the current state of knowledge. Later, you will expand the technical background to become a separate chapter in the half-time report and the final thesis.
\end{instructions}
\end{onlyplan}

\begin{notfinal}
\section{UtExpo Display\sv{UtExpo}}
\label{sec:utexpo}
\begin{instructions}
\begin{onlyplan}
In your project plan, describe what you plan to show at UtExpo. Is it a circuit that does something specific or a software that does something specific? What makes your exhibition at UtExpo interesting?

Here, you can see how to introduce an image in a figure. The picture shown in Figure~\ref{fig:lion-plan} is merely for decorative purposes and shows a cartoon of a lion winking while handling books.

\begin{figure}[htbp]  % order of priority: h here, t top, b bottom, p page
  \centering
  \includegraphics[width=.5\textwidth]{figures/lion.png}
  \caption[The \LaTeX{} logotype]{The \LaTeX{} logotype}
  \label{fig:lion-plan}
\end{figure}

You can refer to it as Figure~\ref{fig:lion-plan} and in this template it is very likely that it ends up on the \emph{next page}. \emph{Please observe:} \LaTeX{} has very good reasons to put figures where they are put after you compile your work. Don't fight it, it will not be worth the effort. Instead, be very specific in your text what the figure describes and refer to it properly.
\end{onlyplan}
\begin{onlymidterm}
Briefly update the reader on your plans for UtExpo. What will you show, has it changed since the Project Plan, and what remains to be done before the exhibition?
\end{onlymidterm}
\end{instructions}
\end{notfinal}
